\subsection{Finite-Dimensional Riesz Representation}

\begin{theorem} \label{thm:riesz-representation-finite-dimensional} \hlt{Riesz Representation Theorem}\\
Let $\mathcal{V}$ be a finite-dimensional inner-product space and let $\varphi:\mathcal{V}\to\mathbb{F}$ be a linear functional. Then there is a unique vector $\mathbf{u}\in\mathcal{V}$ such that $\varphi(\mathbf{v})=\innerprod{\mathbf{v},\mathbf{u}}$ for every $\mathbf{v}\in\mathcal{V}$.
\begin{proof}
Let $(\mathbf{e}_1,\ldots,\mathbf{e}_n)$ be an orthonormal basis of $\mathcal{V}$.\\
By Theorem \ref{thm:orthonormal-basis-coordinates}, $\varphi(\mathbf{v}) =\varphi\left(\sum\limits_{j=1}^{n}\innerprod{\mathbf{v},\mathbf{e}_j}\mathbf{e}_j\right) =\sum\limits_{j=1}^{n}\innerprod{\mathbf{v},\mathbf{e}_j}\varphi(\mathbf{e}_j)$.\\
Set $\mathbf{u}=\sum\limits_{j=1}^{n}\overline{\varphi(\mathbf{e}_j)}\,\mathbf{e}_j$. Since inner product is conjugate linear in second slot, $\innerprod{\mathbf{v},\mathbf{u}} =\sum\limits_{j=1}^{n}\innerprod{\mathbf{v},\mathbf{e}_j}\varphi(\mathbf{e}_j) =\varphi(\mathbf{v})$.\\
For uniqueness, suppose $\innerprod{\mathbf{v},\mathbf{u}_1}=\innerprod{\mathbf{v},\mathbf{u}_2} \ \forall \mathbf{v}\in\mathcal{V}$. Taking $\mathbf{v}=\mathbf{u}_1-\mathbf{u}_2$ gives $\norm{\mathbf{u}_1-\mathbf{u}_2}^2=0$, hence $\mathbf{u}_1=\mathbf{u}_2$.\\
\hlt{Structural Commentary}: An orthonormal basis packages a functional as one representing vector.\\
\hlt{Key Insight}: Finite-dimensional functionals are inner products with fixed vectors.
\end{proof}
\end{theorem}

\begin{definition} \label{def:riesz-vector} \hlt{Riesz Vector}\\
Let $\mathcal{V}$ be finite-dimensional and let $\varphi:\mathcal{V}\to\mathbb{F}$ be a linear functional. The unique vector $\mathbf{u}\in\mathcal{V}$ satisfying $\varphi(\mathbf{v})=\innerprod{\mathbf{v},\mathbf{u}}$ for every $\mathbf{v}\in\mathcal{V}$ is the \hlt{Riesz vector} representing $\varphi$.
\end{definition}

\begin{theorem} \label{thm:riesz-map-conjugate-isomorphism} \hlt{Riesz Map}\\
Let $\mathcal{V}$ be finite-dimensional inner-product space, $\tau:\mathcal{V}\to\mathcal{V}'$, $\tau(\mathbf{u})=\innerprod{\cdot,\mathbf{u}}$. Then $\tau$ is bijective, conjugate-linear.\\
Hence $R=\tau^{-1}$ is conjugate-linear, $R(a\varphi+\psi)=\overline{a}\,R(\varphi)+R(\psi)$, $a\in\mathbb{F},\ \varphi,\psi\in\mathcal{V}'$.\\
Over $\mathbb{R}$ this is ordinary linearity.
\begin{proof}
For $a\in\mathbb{F}$ and $\mathbf{u},\mathbf{w}\in\mathcal{V}$, $\tau(a\mathbf{u}+\mathbf{w}) =\overline{a}\tau(\mathbf{u})+\tau(\mathbf{w})$, so $\tau$ is conjugate-linear.\\
If $\tau(\mathbf{u})=0$, then $\norm{\mathbf{u}}^2=\tau(\mathbf{u})(\mathbf{u})=0$, so $\tau$ is injective. Surjectivity is Theorem \ref{thm:riesz-representation-finite-dimensional}.\\
For $a\in\mathbb{F}$, $\varphi,\psi\in\mathcal{V}'$, and $\mathbf{v}\in\mathcal{V}$, $\innerprod{\mathbf{v},\overline{a}R(\varphi)+R(\psi)} =a\innerprod{\mathbf{v},R(\varphi)}+\innerprod{\mathbf{v},R(\psi)} =(a\varphi+\psi)(\mathbf{v})$.\\
Uniqueness gives $R(a\varphi+\psi)=\overline{a}R(\varphi)+R(\psi)$.\\
\hlt{Structural Commentary}: Riesz identifies $\mathcal{V}$ with $\mathcal{V}'$ conjugate-linearly.\\
\hlt{Key Insight}: The representing slot conjugates scalars.
\end{proof}
\end{theorem}

\begin{figure}[H]
  \centering
  \begin{tikzpicture}[
  >=Stealth,
  every node/.style={font=\small},
  space/.style={draw=gray!45, rounded corners=3pt, minimum width=3.7cm, minimum height=2.5cm},
  dual/.style={draw=gray!45, rounded corners=3pt, minimum width=3.7cm, minimum height=2.5cm, fill=gray!5},
  vec/.style={->, blue!75!black, line width=0.9pt},
  map/.style={->, gray!70, line width=0.75pt},
  test/.style={->, orange!80!black, line width=0.7pt}
]
  \node[space] (V) at (0,0) {};
  \node[anchor=north west] at (-1.65,1.1) {$\mathcal{V}$};
  \coordinate (O) at (-0.95,-0.65);
  \coordinate (u) at (0.75,0.65);
  \coordinate (v) at (0.8,-0.15);
  \draw[vec] (O) -- (u) node[midway, above left] {$\mathbf{u}$};
  \draw[test] (O) -- (v) node[midway, below] {$\mathbf{v}$};

  \node[dual] (Vdual) at (5.3,0) {};
  \node[anchor=north west] at (3.65,1.1) {$\mathcal{V}'$};
  \node[align=center] at (5.3,0.25) {$\varphi_{\mathbf{u}}$};
  \node[align=center] at (5.3,-0.45) {$\varphi_{\mathbf{u}}(\mathbf{v})=\innerprod{\mathbf{v},\mathbf{u}}$};

  \draw[map] (1.95,0.25) -- (3.35,0.25) node[midway, above] {Riesz};
  \draw[test] (1.0,-0.15) .. controls (2.5,-1.55) and (4.0,-1.55) .. (5.0,-0.85)
    node[midway, below] {evaluation};
\end{tikzpicture}

  \caption{Riesz identifies a functional with its representing vector.}
  \label{fig:riesz-representation-vector-functional}
\end{figure}

\begin{perspective} \label{per:riesz-representation-coordinates} \hlt{Riesz as Coordinate Duality}\\
The vector $\sum_{j=1}^{n}\overline{\varphi(\mathbf{e}_j)}\mathbf{e}_j$ seems to depend on the orthonormal basis, but Theorem \ref{thm:riesz-representation-finite-dimensional} shows that the resulting representing vector is intrinsic.
\end{perspective}
